% =============================================================================
% CHAPTER 3: Methodology (or ``Specification, Design and Implementation'' / ``Approach'')
% =============================================================================
% SPDX-FileCopyrightText: 2026 Frank Langbein <frank@langbein.org>
% SPDX-License-Identifier: LPPL-1.3c
%
% This chapter describes how your solution works. Title it Methodology, Approach,
% or Specification, Design and Implementation; adapt or split as needed. Not all
% sections below apply to every project---pick, omit, or rename.
%
% Suggested sections by project type:
%   System design:  Specification, System design, Implementation. Optional: Datasets
%                   if data is central; Ethical considerations.
%   Research:       Research design, Literature and evidence (search strategy,
%                   inclusion criteria), Specification (= problem + requirements),
%                   Approach and solution, Implementation or experimental setup.
%   AI/ML:          Specification, Datasets, AI/ML models, System design (if a
%                   larger system), Implementation. Optional: Research design,
%                   Literature and evidence; Summary.
% =============================================================================

This chapter describes how the problem was addressed. Choose sections that fit your project type; omit or merge as needed. Replace with enough detail for reproduction.


\section{Research design}\label{sec:research-design}
% Research: overall approach, methods (e.g.\@ quantitative, qualitative, design science), rationale. Omit for pure system-design projects.


\section{Specification}\label{sec:spec}
% System design: requirements (what the system must do; optional subsections: functional, non-functional). Research: problem and what is required of a solution.


\section{Datasets}\label{sec:datasets}
% AI/ML / empirical: data sources, preprocessing, train/validation/test split. Omit if not data-driven.


\section{AI and machine learning models}\label{sec:aiml}
% AI/ML: model choice, training and validation, evaluation methodology. Omit otherwise.


\section{System design}\label{sec:design}
% System design: architecture, data flow, behaviour (UML, ER, state charts). Use diagrams. Research: use ``Approach and solution'' for your proposed solution(s).


\section{Implementation}\label{sec:impl}

\subsection{Introduction}

The implementation of this system utilises ephermeral Linux containers, injecting intentional vulnerabilites into them, exposing them to a student through a controlled network path, verifiying the challenge through a flag, and prov

\subsection{Overall System Architecture}
\subsection{Infrastructure orchestration and ephemeral lab provisioning}
\subsubsection{Proxmox and container model}
\subsubsection{Ansible as the orchestration layer}
\subsubsection{Network access and learner connectivity}
\subsubsection{Automatic cleanup and the TTL reaper}
\subsection{LLM driven vulnerability generation and injection}
\subsection{Web application}




\section{Summary}\label{sec:summary}
% Optional: short summary of the methodology. Omit if the chapter is already concise.
